\scanpage{27}
% archive sheet number: 26
% Upper argument cancelled by long vertical strokes.
% Source lines 1--4.
\deleted{D'autre part, \unsure{le degré} des $K_s$ ($s\in S$) étant fonction constructible, \unsure{ne prend qu'un|on peut prendre} \unsure{nb} fini de valeurs, soit $n$ un multiple commun. Considérons}
\diagram{\begin{tikzcd}[ampersand replacement=\&]
\deleted{\unsure{K|G}}\arrow[r,"{\deleted{\unsure{u_n|u}=n.\operatorname{id}_{K}}}"]\&\deleted{K}
\end{tikzcd}}
% Source line 5, locally struck as well as block-cancelled.
\deleted{\deleted{\unsure{C'est une isogénie} $\unsure{N}$ \unsure{est alors nul}}}
% Source lines 6--11.
\deleted{Alors ${}_nK=\operatorname{Ker}u_n$, or \uline{la section unité de $K$ est une immersion ouverte} donc ${}_nK$ est un ss-schéma \unsure{ouvert fermé} de $K$. Comme il contient les fibres de $K$ par construction, il est égal à $K$ \unsure{ensemblist\textsuperscript{t}} donc aussi \unsure{schématiqu\textsuperscript{t}|schématiquement}.}

% Lower NB cancelled by long diagonal/vertical strokes.
% Source lines 12--14.
\deleted{\uline{NB} La condition \inserted{``étale''} devient fausse si on suppose seulement $G$ loc. de présentation finie sur $S$,}
\margin{\deleted{ex : \unsure{schémas de Néron infini.|schémas de Néron fini.|schémas de Néron infinis.|schémas de Néron finis.}}}
% margin note written sideways
\margin{\deleted{\deleted{\unsure{Soit}} D'autre part}}
% Source lines 14--19.
\deleted{Une application intéressante : Supposons pour $X/S$ plat propre, \deleted{\unsure{NB}} $\underline{\operatorname{Pic}}^{\unsure{\tau|t}}_{X/S}$ \deleted{\unsure{et $\underline{\operatorname{Pic}}$}} existant, \uline{supposons} que en un pt $\unsure{s|s_0}$ $\underline{\operatorname{Pic}}^0_{\unsure{X_s|X_{s_0}|X_0}/\unsure{k(s)|\kappa(s)}}$ est \inserted{\unsure{abélien}} simple $/\unsure{k(s)|\kappa(s)}$, alors $\underline{\operatorname{Pic}}^0_{X/S}$ existe au voisinage de $s_0$.}
% Source lines 20--24.
\deleted{De plus, $\underline{\operatorname{Pic}}^{\tau}_{X/S}/\underline{\operatorname{Pic}}^0_{X/S}$ existe, et est \unsure{non ramifié} sur $S$. Il est même \uline{étale} sur $S$. Donc [si $\underline{\operatorname{Pic}}^{\tau}_{X/S}$ est propre sur $S$] on a la constance de la torsion de N.S.}
